\section{Representation algebra and declared cost models}
\subsection{Why a count theorem can have a short proof}
\code{PacketCostModels} defines finite address and event types before counting them. The block widths are natural numbers. A diagonal block has a pair of within-block indices. A thin factor has one ambient and one rank index. An application event additionally has a right-hand-side index and a Boolean distinguishing multiplication from accumulation. The Boolean gives the declared two-scalar-operation convention.
\excerpt{}{bdlr_count}
\walk{The type \code{BDLROperation} is a sum of the diagonal-action, right-factor, and left-factor event types. \code{simp} evaluates cardinalities of finite products and sums, and \code{ring} rearranges the resulting polynomial. The count includes the accumulation into the diagonal result through the declared third kernel. It does not establish how many GPU instructions run.}

The storage count concerns stored scalar values. Sparse index metadata and workspace are outside the corresponding models. The materialization theorem counts an explicit dense low-rank multiplication followed by one addition per output entry. Its quadratic remainder is exactly \(n^2\). This is a count for that declared construction, rather than a lower bound for every conceivable implementation.
\excerpt{}{band_count}
\walk{The banded-count development counts the actual subtype of pairs in \code{Fin p × Fin p} whose distance is at most \code{w}. The premise \code{w < p} controls the two triangular boundary losses. The preceding recurrence and finite-sum lemmas establish the cardinality. This endpoint presents the closed formula. Replacing \(p\) by \(n/b\) additionally requires positive block width and exact divisibility.}

\subsection{Capacity inequalities are hypotheses}
\excerpt{}{roofline}
\walk{Positive peak arithmetic and bandwidth rates make division preserve order. The assumed inequalities \code{F ≤ P * T} and \code{B ≤ W * T} yield the two lower bounds on time, then \code{max_le} combines them. Lean does not measure a device rate or validate that an experimental execution satisfies these capacity assumptions.}

\subsection{An exact low-rank product}
\contract{\code{LowRankAlgebra} works with actual finite matrices over a field with an involution. The left and right factors have explicit possibly different width types.}
\excerpt{}{bdlr_product}
\walk{The constructors \code{bdlrProductLeft} and \code{bdlrProductRight} concatenate columns, so the new correction width is the sum of the two input widths. The proof expands the concatenated product, reverses conjugate transposes of products, distributes multiplication, and normalizes additive terms. No rank identity is assumed in this algebraic equality. A separate theorem bounds the correction rank by its column count. Packet diagonality of the base product is also proved separately.}
\excerpt{}{chain_factors}
\walk{The empty chain is the identity and has an empty correction-column type. For the successor step, rewrite the current chain using the induction hypothesis and invoke the actual two-factor identity. The source thereby constructs factors for every length. The subsequent rank theorem bounds their width by the sum of input widths, including zero widths.}

\subsection{Woodbury exposes both invertibility assumptions}
\excerpt{}{woodbury}
\walk{\code{hD} proves the base matrix is a unit. \code{hS} proves the small update matrix is a unit. The endpoint specializes Mathlib's Woodbury identity, with the middle rank-space matrix equal to one. The update's invertibility, batch solve formula, and L2 operator norm bound are independent declarations in the same module. A singular base or singular update matrix cannot be passed through the displayed theorem simply because an inverse notation exists.}
